\begin{definition}
Функция $f$ \textit{конструируема по времени}, если существует МТ, которая на входе $1^n$ вычисляет $f(n)$ за $O(f(n))$ 
\end{definition}
\begin{theorem}\label{th:hir-time}
Пусть $f, g$ --- строго монотонные конструируемые по времени функции и 
 \\$f(n)\log{f(n)} = o(g(n))$. Тогда $\dtime{f(n)} \subsetneq \dtime{g(n)}$
\end{theorem}
\begin{proof}
\begin{enumerate}
    \item Построим функцию $h$ -- конструируемую по времени, такую что $h(n)\log{h(n)} = O(g(n))$ и $f(n) = o(h(n))$. Например, сгодится $\frac{g(n)}{\log{g(n)}}$
     \item Рассмотрим множество $D = \{M: M(M)=0$ и останавливается за $h(|M|)$ шагов$\}$

     \item Утверждение (б/д) $D \in \dtime{h(n)\log{h(n)}}$
     \item Из предыдущего пункта следует, что $D \in \dtime{g(n)}$, так как мы выбирали $h$ такую, что $h(n)\log{h(n)} = O(g(n)))$.
     \item Предположим, что $D \in \dtime{f(n)}$ и распознается машиной Q. 
     $$Q(M) = 1 \longleftrightarrow M \in D$$ по определению, значит 
     $$Q(Q) = 1  \longleftrightarrow Q \in D$$
     Значит, если $Q \in D$, то $Q(Q) = 0$ и машина останавливается за $h(|Q|)$ шагов, что сразу приводит к противоречию. Если же $Q \not\in D$, то $Q(Q) = 0$, то есть у нас не выполняется условие, что она останавливается за $h(|Q|)$ шагов. Можно раздуть, дописывая в конец фиктивные символы так что $C\cdot f(|Q|) < h(|Q|)$, так как $f(n) = o(h(n))$, и снова придем к противоречию, так как $D \in \dtime{f(n)}$, значит, распознается за $O(f(|Q|))$\end{enumerate}
\end{proof}